\pagestyle{plain}

\section{À propos de cette publication}



\subsection{Pourquoi les objets connectés?}


Alors que dans certaines disciplines le temps commence à manquer pour traiter l'ensemble du programme, certains évoquent déjà l'idée d'en faire plus !

En effet, les enseignants utilisent déjà les outils numériques. Par exemple, dans les classes de mathématiques, l'utilité du tableur et de GeoGebra n'est plus à démontrer. Jusqu'à l'introduction de l'algorithmique, ces deux logiciels efficaces et maîtrisés par les enseignants étaient amplement suffisants.
Est-ce donc juste un effet de mode de faire cours avec les robots (Thymio, Mbot), les objets  programmables et connectés (Arduino, \mb, STM education, Raspberry Pi) ou est-ce une nouvelle façon d'aborder notre enseignement?
Ces nouvelles possibilités technologiques, forcément chronophages, nous permettront-elles de traiter un contenu disciplinaire exigeant dans un cadre institutionnel contraignant?

Nous n'avons bien sûr pas toutes les réponses à ces questions mais nous pensons que lorsqu'il est accompagné de certains de ces outils, notre enseignement a beaucoup à y gagner.

L'introduction de l'algorithmique en lycée professionnel nous interroge. Longtemps il nous a semblé impensable et inenvisageable d'avoir à enseigner un langage de programmation comme Python auprès d'un public d'élèves globalement en difficulté avec les mathématiques. Fort de ce constat, nous avons cherché les moyens de lier les mathématiques à la logique et au raisonnement algorithmique. C'est pourquoi nous avons exploré les potentialités des objets connectés.


\begin{formule}
	Notre postulat est double. Nous pensons que :
		\begin{itemize}
			\item grâce à des situations réelles et concrètes, les objets connectés facilitent la mise en activité de tous les élèves ;
			\item grâce à des activités simples mais évolutives centrées autour de réalisations matérielles, la dimension affective du travail est valorisée.  Soyons fous et espérons que l'élève tisse une histoire personnelle avec l'activité, qu'il soit fier du travail accompli et  qu'il prenne également du plaisir à expliquer et à montrer ses réalisations.
	\end{itemize}
\end{formule}


En devenant de plus en plus simples, accessibles et facilement utilisables, les objets connectés permettent d'aborder des contenus disciplinaires et de développer des compétences transversales essentielles pour l'élève.

En travaillant à partir des objets connectés, la situation de départ est plus concrète et l'objectif à atteindre suffisamment clair pour l'élève. Plus ou moins guidé selon son niveau d'expertise technique, il est alors libre dans sa démarche.
Avec des interfaces de programmation accompagnées parfois de simulateurs, la démarche par essais et erreurs a ici toute sa place.
Par ailleurs, l'élève devra clarifier sa pensée avant de verbaliser ses idées en langage naturel. Il pourra ainsi proposer et élaborer un modèle acceptable par la machine pour enfin traduire son algorithme en se pliant à la rigueur du langage de programmation. 

Effectuant régulièrement des va-et-vient entre abstraction et réalité, cherchant à valider son algorithme à partir d'un visuel ou d'une exploitation des résultats, l'élève entre progressivement dans la modélisation.

Les scénarios proposés dans cette brochure permettent tout cela : une approche des mathématiques et des sciences qui laisse la place à l'expérimentation : manipulation, programmation et auto-validation. 




\pagebreak

\subsection{Qui sommes-nous?}

Nous sommes des enseignants de maths/sciences regroupés au sein d’un groupe de recherche de l'IREM de Marseille.

\begin{center}
	\href{http://url.univ-irem.fr/ineflp}
	{\includegraphics[width=0.5\linewidth]{res/fig-logo-ineflp.png}}
\end{center}

Notre groupe, Innovation, Expérimentation et Formation en Lycée Professionnel (InEFLP) consacre une partie de son travail à l’enseignement de l’algorithmique en classes de lycée professionnel. Dans le cadre de cette recherche, nous explorons les objets connectés tels que Arduino, \mb, \st ou mbot.





\subsection{Liens utiles}
\begin{description}
    \item[Page du groupe InEFLP]  ~\\
        \url{http://url.univ-irem.fr/ineflp}
    \item[IREM de Marseille] Site académique de l'IREM de Marseille\\
        \url{http://url.univ-irem.fr/mars}
    \item[Portail des IREM] Site national des IREM\\
        \url{http://www.univ-irem.fr/}
    \item[Formation à l'algorithmique LP et SEGPA]
        Padlet de utilisé lors de nos formations académiques\\
        \url{http://url.univ-irem.fr/stage-algo}
    \item[Collecte de ressources pour \mb]
        Padlet sur \mb utilisé en formation\\
        \url{http://url.univ-irem.fr/algo2017-microbit}
    \item[Brochure sur \mb]
        Publication de la \href{http://www.univ-irem.fr/spip.php?rubrique18}{C2i TICE} pour une prise en main de \mb\\
        \url{http://url.univ-irem.fr/c2it-mb-t1-pdf}
    \item[Description \mb] Fiche sommaire de description de \mb\\
        \url{http://url.univ-irem.fr/ineflp-microbit}
    \item[Site IREM dédié à \mb] Site de ressources sur \mb du groupe\\
        \url{http://url.univ-irem.fr/o}
\end{description}


\begin{center}
    \href{http://url.univ-irem.fr/ineflp}
    {\includegraphics[width=0.1\linewidth]{res/fig-ineflp-qr.png}}
\end{center}

